\chapter{Conclusiones}\label{chapter:conclusiones}

Este trabajo ha diseñado, implementado y evaluado un sistema \englishText{end-to-end} que transforma
dictados médicos en informes clínicos estructurados e interoperables, integrando captura segura de
audio, transcripción \englishText{ASR}, extracción y etiquetado (\englishText{NER}) y normalización de entidades clínicas
(\englishText{SNOMED CT}/\englishText{CIE-10-ES}), generación de informes mediante \englishText{LLMs}
y exportación en formato \englishText{HL7 FHIR}. El sistema ha sido construido y desplegado de forma
funcional, y evaluado tanto a nivel de módulo individual como de \textit{pipeline} completo, lo que
permite extraer conclusiones tanto sobre su viabilidad técnica como sobre sus limitaciones actuales
de cara a un uso clínico real.

\section{Conclusiones generales}\label{section:conclusiones_generales}

El resultado más relevante de la evaluación, discutido en la
Sección~\ref{section:discusion_resultados}, es que el desempeño de un sistema clínico de este tipo no
puede inferirse a partir del desempeño aislado de sus componentes. Cada módulo evaluado de forma
independiente frente a su propio \textit{gold standard} obtiene resultados notablemente sólidos: un
\englishText{WER} medio de $\approx$6,5\,\% en \englishText{ASR}, un 85,5\,\% de acierto micro en la
normalización a \englishText{SNOMED CT} y un F1 micro de hasta 43,8\,\% en \englishText{NER} bajo
criterio granular. Sin embargo, al integrarse en un \textit{pipeline} secuencial, donde cada etapa
consume la salida de la anterior sin posibilidad de corrección aguas abajo, estos mismos componentes
muestran un rendimiento muy inferior en la evaluación \englishText{end-to-end} (puntuaciones de
1,55 a 2,90 sobre 5), y únicamente entre el 20\,\% y el 30\,\% de los informes generados resultan
aceptables para el veredicto del \englishText{LLM-as-judge}. Esto confirma que la arquitectura
secuencial propaga y amplifica errores moderados de una etapa (una sigla mal transcrita, una entidad
truncada por sobre-segmentación) hasta convertirlos en errores clínicos graves en el informe final
(diagnósticos sustituidos, dosis erróneas, inversiones de hallazgos), un comportamiento de
acumulación de error que constituye el hallazgo central de este trabajo y que cualquier
arquitectura similar basada en encadenamiento de modelos debe asumir y mitigar explícitamente.

En cuanto al cumplimiento de los objetivos específicos planteados en la
Sección~\ref{section:objetivos-especificos} (\nameref{section:objetivos}, Sección
\ref{section:objetivos}), los resultados son dispares. El objetivo~1 (\englishText{WER}
$\le$10--12\,\%) se cumple con margen, demostrando que los modelos \englishText{ASR} actuales,
adaptados al léxico médico en español, son suficientemente maduros para esta tarea en entornos
controlados. El objetivo~4 (validación estructural \englishText{FHIR} $>$95\,\%) se cumple en
términos estrictamente sintácticos --- el 100\,\% de los \textit{Bundles} generados son
estructuralmente válidos ---, si bien con la reserva importante de que dicha validación no garantiza
la corrección semántica del contenido clínico que representa. Por el contrario, los objetivos~2
($F_1 \ge 0,85$ en \englishText{NER}) y~3 (aceptación médica $\ge$80\,\%) no se cumplen, con
resultados muy por debajo de los umbrales fijados. El objetivo~2 falla principalmente por
sobre-segmentación de \textit{spans} y por la falta de reconocimiento de siglas y abreviaturas
clínicas, mientras que el incumplimiento del objetivo~3 es, en gran medida, una consecuencia directa
de la propagación de errores ya descrita: el informe \englishText{LLM} hereda y amplifica los fallos
de las etapas \englishText{NER} y de normalización terminológica que lo preceden.

Estos resultados permiten concluir que el sistema desarrollado constituye una prueba de concepto
técnicamente viable y arquitectónicamente completa --- cubre de extremo a extremo el flujo desde el
audio hasta un recurso \englishText{FHIR} interoperable, respetando los requisitos de privacidad y
seguridad planteados en la motivación del trabajo ---, pero que, en el estado actual, no alcanza el
nivel de fiabilidad clínica necesario funcionar sin supervisión humana. En definitiva, el trabajo demuestra
que la dificultad de este tipo de sistemas no reside
únicamente en alcanzar buen desempeño en cada tarea de \englishText{PLN} por separado, sino en
diseñar mecanismos de robustez y de propagación controlada del error a lo largo de todo el
\textit{pipeline}.

\section{Trabajos futuros}\label{section:trabajos_futuros}

A partir de las limitaciones identificadas en la Sección~\ref{section:discusion_resultados}, se
proponen las siguientes líneas de continuación:

\begin{itemize}
    \item \textbf{Mitigación de la propagación de errores entre etapas:} introducir mecanismos de
    verificación intermedios entre módulos (por ejemplo, validación de la consistencia entre las
    entidades \englishText{NER} y el texto fuente antes de pasarlas al módulo de normalización, o
    confirmación cruzada de las negaciones detectadas) en lugar de un encadenamiento puramente
    secuencial sin retroalimentación.

    \item \textbf{Mejora del módulo \englishText{NER} sobre los patrones de error identificados:}
    incorporar un diccionario o módulo de expansión de siglas y abreviaturas clínicas, y ajustar la
    delimitación de \textit{spans} para reducir la sobre-segmentación de entidades compuestas
    (diagnósticos en varias palabras, valores numéricos con unidades), que concentran la mayor parte
    de los errores observados.

    \item \textbf{Normalización terminológica sensible al contexto:} sustituir o complementar la
    búsqueda por coincidencia de texto sobre la entidad aislada por un enlace de entidades
    (\textit{entity linking}) que tenga en cuenta el contexto clínico del dictado completo, lo que
    debería mejorar la precisión de la baja cobertura observada en \englishText{CIE-10-ES}.

    \item \textbf{Reducción de alucinaciones del informe \englishText{LLM}:} explorar técnicas de
    generación aumentada por recuperación (\englishText{RAG}) o de verificación post-generación que
    contrasten cada afirmación del informe contra las entidades y el texto fuente, para detectar
    inversiones de hechos clínicos, errores de magnitud en dosis y expansión incorrecta de siglas
    antes de aceptar el informe como definitivo.

    \item \textbf{Evaluación sobre infraestructura de mayor capacidad:} replicar la evaluación
    \englishText{end-to-end} sobre \textit{hardware} con mayor VRAM disponible, lo que permitiría
    emplear variantes de \englishText{LLM} de mayor escala que \texttt{llama3.1:8b} y determinar en
    qué medida las limitaciones de calidad observadas (alucinaciones, colapsos de contenido) se deben
    al tamaño del modelo o son inherentes a la arquitectura secuencial del \textit{pipeline}.
\end{itemize}
